        These attributes are not case-sensitive but will be converted to
        lower case.

        The following attributes are used by {\bf input()} itself.
        
        \btabularx
        \hline
        \multicolumn{4}{|l|}{\bf input mode attributes} \\
        \hline
        differential & netlist & log          & differential                         \\
        loc          & netlist & str or []str & locations                            \\
        \hline
        \btabularxend
                
        Where an attribute value is an array the dimension must be equal to
        the number of pins and each array member will apply to the corresponding
        output pad or buffer. A single attribute value will be applied to all pins.

        {\bf differential} indicates that the input pads are differential.

        The {\bf loc} attribute is mandatory. The array dimension
        must exactly match the number of bits in the variable type. If the input
        is differential the array dimension must be twice the number of
        bits, the locations being used in pairs for the +ve and the -ve inputs
        (in that order). The location array matches the associated signal
        in order LSB to MSB for the array members and the signal bits.
            
        Pin strings are matched left to right to an input variable
        starting at the the least-significant bit. For a compound type pins are
        matched left to right from inner to outer nested types. For a struct the
        strings are matched left to right in field order.

        FPGA package pins are usually already defined in a platform header file,
        for example -
        \begin{lstlisting}[gobble=12, language=threepl]
           iopins(a_pins, "[12]str",
                        {
                            "A0", "B0", "A1", "B1", "A2", "C2",
                            "A3", "B3", "A4", "B4", "C3", "C4"
                        }
                );
        \end{lstlisting}
        There might be additional attributes included for these pins.

        \begin{lstlisting}[gobble=12, language=threepl]
           input(v, "[2][2]uint:2", loc=a_pins[2..9]);
        \end{lstlisting}

        The resulting input connections would be -\\
        pins "A1", "B1" (LSB, MSB) are connected to v[0][0] \\
        pins "A2", "C2" (LSB, MSB) are connected to v[0][1] \\
        pins "A3", "B3" (LSB, MSB) are connected to v[1][0] \\
        pins "A4", "B4" (LSB, MSB) are connected to v[1][1] \\

        The following attribute is read-only and is created when {\bf
        input()} is called and can be retrieved from the input mode
        variable if required by using inbuilt function {\bf attribute()}.
        
        \btabularx
        \hline
        \multicolumn{4}{|l|}{\bf input mode} \\
        \hline
        ids          & 3PL & []str & element IDs       \\
        \hline
        \btabularxend       
        
        {\bf ids} is a string array that contains the netlist identifiers of the
        input buffer elements. This attribute is read-only and is initialised
        when the input variable is created. This identifier can be used should
        there be a necessity to apply a constraint (other than location which is
        done via the port name) to the input buffer.

        The following attributes are required for constraint generation and are
        handled at completion of immediate execution by library procedure {\bf
        generateconstraints()} \autorefph{sec:lpgenerateconstraints}. The
        attribute values are single values which will be applied to all input
        bits.
        
        \btabularx
        \hline
        \multicolumn{4}{|l|}{\bf input mode} \\
        \hline
        ibuf\_delay\_value & constraint & int or []int & extra delay to IBUF                      \\
        ifd\_delay\_value  & constraint & int or []int & extra delay to IFD                       \\
        tnm                & constraint & str or []str & timing group id                          \\
        pullup             & constraint & log or []log & pullup                                   \\
        pulldown           & constraint & log or []log & pulldown                                 \\
        diff\_term         & constraint & log or []log & differential  $100\Omega$ termination    \\
        iostandard         & constraint & str or []str & voltage standard                         \\
        \hline
        \btabularxend
        
        {\bf ibufdelay} specifies a delay to be inserted in the non-registered
        output path from an input buffer. The values (integers) are 0 to 12
        inclusive (default 0).
        
        {\bf ifddelay} specifies a delay to be inserted in the registered
        path within an input buffer. The values (integers) are 0 to 6
        inclusive (default is a value generated automatically, see device
        databook).
        
        {\bf tnm} is the name of a timing group in which the input pads
        are to be included for timing constraint purposes.
        
        If {\bf pullup} or {\bf pulldown} are true the input will be given
        a weak pullup or pulldown. It is not allowed for both these
        attributes to be true.
        
        {\bf diff\_term} indicates that a $100\Omega$ termination resistance
        be applied to a differential input.
        
        {\bf iostandard} is a string specifying one of the allowed
        I/O standards (see Xilinx documentation).
        
        The following attributes are tolerated but ignored. They may be included
        as attributes for the output side of a bi-directional pin.
        
        \btabularx
        \hline
        \multicolumn{4}{|l|}{\bf input mode} \\
        \hline
        slew              & constraint & str or []str   & timing group id                          \\
        drive             & constraint & uint or []uint & drive current                                  \\
        \hline
        \btabularxend
        
        Any attributes that are not listed above will result in a fatal error.
